O trabalho demonstrou, em bancada, que uma ponte CAN sem fio baseada em
ESP-NOW pode estender o sistema de aquisição de dados sem abandonar a organização
lógica por identificadores CAN. Essa validação não resolve todos os aspectos de
uso embarcado, mas confirma que a abordagem é tecnicamente plausível para pontos
de instrumentação em que o cabeamento se torna uma limitação prática.

A principal leitura técnica é que a viabilidade do WCAN depende da combinação
entre agrupamento temporal e política de transporte. O agrupamento foi essencial
para reduzir a pressão sobre o rádio, enquanto o \emph{multicast} se mostrou a
alternativa mais favorável para múltiplos fluxos por remover parte da
coordenação do caminho crítico. Dessa forma, o resultado do trabalho não é apenas
a implementação de um enlace sem fio, mas a identificação das condições em que
esse enlace passa a operar com margem útil para aquisição de dados.

As dificuldades observadas delimitam o alcance da solução. O envio amostra a
amostra expôs rapidamente os limites do caminho ESP-NOW, e a confirmação no
nível MAC adotada no \emph{multicast} não equivale à confirmação de
processamento pela aplicação remota. Por isso, o WCAN deve ser entendido como
ponte de instrumentação, não como substituto do CAN físico em funções críticas
do veículo, nas quais continuam relevantes as garantias elétricas, temporais e
de arbitragem do barramento cabeado.

Como continuidade natural, o projeto deve avançar da validação em bancada para
ensaios embarcados no veículo, com integração ao barramento CAN em condição real
e avaliação da robustez operacional do nó físico em uso prolongado. Essa etapa
permitirá verificar, no ambiente final de aplicação, os efeitos de montagem,
alimentação, posicionamento das antenas e interferências sobre a estabilidade da
aquisição.
